\section{The proved local bound at low height}
\label{sec:low-cubic}
%======================================================================

This is the closed part of the bridge.  Geometry alone bounds the low-height increment in each nontrivial square block, and the isolated unit source is a finite endpoint.  The resulting cumulative low sector has translated-window energy of order $HN^2$.

Fix $\Lambda>0$.  Use the canonical conventions
\[
 \Pplus(1)=1,\qquad c_1=1,\qquad Y_1=0,
\]
and assign every integer $m$ to its square block
\begin{equation}
 j=e(m)=\lfloor\sqrt m\rfloor.
 \label{eq:entry-block}
\end{equation}
Define
\begin{align}
 d_j^{\mathrm{low}}(\Lambda)
 &:={}
 \sum_{\substack{m\in B_j\\|Y_m|\le\Lambda j}}\mu(m),
 \label{eq:low-block-increment}\\
 d_j^{\mathrm{high}}(\Lambda)
 &:={}
 \sum_{\substack{m\in B_j\\|Y_m|>\Lambda j}}\mu(m).
 \label{eq:high-block-increment}
\end{align}
The predicates are complementary, so the complete block increment satisfies
\begin{equation}
 \sum_{m\in B_j}\mu(m)
 =d_j^{\mathrm{low}}(\Lambda)+d_j^{\mathrm{high}}(\Lambda).
 \label{eq:block-low-high-recombination}
\end{equation}

The block $B_0=[0,1)$ contains only $0$, whose M\"obius weight is zero, so
\begin{equation}
 d_0^{\mathrm{low}}(\Lambda)=d_0^{\mathrm{high}}(\Lambda)=0.
 \label{eq:zero-block-endpoint}
\end{equation}
The unit source $m=1$ lies in $B_1=[1,4)$ and is low because $Y_1=0$.  For every block, the nontrivial low sources satisfy the occupancy bound in \cref{eq:constant-occupancy}; hence
\begin{equation}
 \boxed{
 |d_j^{\mathrm{low}}(\Lambda)|
 \le \lfloor\Lambda\rfloor+\1_{\{j=1\}}
 \le \lfloor\Lambda\rfloor+1.
 }
 \label{eq:low-block-bound}
\end{equation}
Equivalently, one may separate the single $m=1$ contribution and retain the sharper bound $\lfloor\Lambda\rfloor$ for the remaining nontrivial population.

Define the cumulative sectors
\begin{align}
 S_n^{\mathrm{low}}(\Lambda)
 &:={}
 \sum_{0\le j\le n}d_j^{\mathrm{low}}(\Lambda),
 \label{eq:low-prefix}\\
 S_n^{\mathrm{high}}(\Lambda)
 &:={}
 \sum_{0\le j\le n}d_j^{\mathrm{high}}(\Lambda).
 \label{eq:high-prefix}
\end{align}
Summing \cref{eq:block-low-high-recombination} over $0\le j\le n$ gives the exact native decomposition
\begin{equation}
 \boxed{
 S_n=M((n+1)^2-1)
 =S_n^{\mathrm{low}}(\Lambda)+S_n^{\mathrm{high}}(\Lambda).
 }
 \label{eq:low-high-split}
\end{equation}

For $n\ge1$, the blocks $B_1,\ldots,B_n$ contain all sources up to $X_n$, the unit source occurs once, and therefore
\begin{equation}
 |S_n^{\mathrm{low}}(\Lambda)|
 \le1+\lfloor\Lambda\rfloor n
 \le(\lfloor\Lambda\rfloor+1)n.
 \label{eq:low-prefix-bound}
\end{equation}
Consequently, if $N\ge1$ and $1\le H\le N$, then every $n\in[N,N+H)$ satisfies $n<2N$, and hence
\begin{equation}
 \boxed{
 \sum_{n=N}^{N+H-1}
 |S_n^{\mathrm{low}}(\Lambda)|^2
 \le
 4(\lfloor\Lambda\rfloor+1)^2HN^2.
 }
 \label{eq:low-local-bound}
\end{equation}
The global cubic consequence is
\begin{equation}
 \sum_{n\le N}|S_n^{\mathrm{low}}(\Lambda)|^2
 \ll_{\Lambda}N^3.
 \label{eq:low-cubic}
\end{equation}
No M\"obius cancellation is used.
\leanxref{Lean proof}{canonical low-height occupancy and low-increment control}{https://github.com/OVVO-Financial/RH_Square_Block/blob/main/lean/RHLean/Analysis/CanonicalLowOccupancy.lean}

Finally, from \cref{eq:low-high-split},
\begin{align}
 |S_n|^2
 &\le2|S_n^{\mathrm{low}}|^2+2|S_n^{\mathrm{high}}|^2,
 \label{eq:total-from-low-high}\\
 |S_n^{\mathrm{high}}|^2
 &\le2|S_n|^2+2|S_n^{\mathrm{low}}|^2.
 \label{eq:high-from-total-low}
\end{align}
These two norm inequalities, rather than an invalid energy subtraction identity, give the exact equivalence between the total and high-sector local criteria.

%======================================================================
